\section{Experimental Setup}
\label{sec:methods}
% src: work/datasets_v2/DESIGN.md (tiers, schema, authoring pipeline); questions/*.jsonl (counts);
%      TODO.md "User review (2026-09-25)" and review files (review outcomes); lex4_replay/REPLAY.md (decisive spans)
% src: paper/analysis/revision_stats.md "Families" (labels 47, constituents 33, components 20, blocks 29)
% src: RAG-Paper-upgrades/ragbench: retrieval/hybrid.py, config.yaml (chunking), approaches/hybrid.py (rerank),
%      approaches/lex5b/{approach,views}.py, approaches/lex8/fusion.py; work/datasets_v2/setr/SETR.md;
%      test_freeze/PROTOCOL.md sec. 2 and 9 (arms, readers, pins, gates); lex5b/LEX5B.md (token amendment)
% src (dev record): lex4_replay/REPLAY.md, lex11/LEX11.md, lex5b/LEX5B.md, devrun/RESULTS_DEV.md; Table negative-full in Appendix
% src (number check 2026-11-11): 48 of 54 = work/datasets_v2/lex8/LEX8.md "Diagnosis" (also lex9/LEX9.md); 32 of 156 (31 good, 1 fix) = work/TODO.md line 95;
%      35 of 4,608 unresolved = TEST_RESULTS.md sec. 6 (35 cells) and S5_stats/RESULTS.md sec. 0; pool chars = work/multidoc/pools/<pool>.json (fa 11.64M, en 26.42M)
% Rewrite 2026-11-10: benchmark, pipelines, readers, judge and protocol merged into one section; protocol history in one paragraph.

Four things define the setup, namely the benchmark, the four pipelines, the readers and the judge, and the rules fixed before the test.

\subsection{Benchmark}
\label{sec:benchmark}
The benchmark pairs natural Persian and English documents with questions of three difficulty tiers.
% src: work/datasets_v2/DESIGN.md (tier definitions, integration operations); questions/*_{dev,test}.jsonl (counts per tier,
%      language and domain); examples verbatim from questions/law_dev.jsonl V2-dev-LAW-en-T1-003,
%      rpt_dev.jsonl V2-dev-RPT-en-T2-003, rpt_dev.jsonl V2-dev-RPT-en-T3-001, rpt_dev.jsonl V2-dev-RPT-en-AB-001 (dev questions)
\begin{table*}[t]
\centering
\small
\renewcommand{\arraystretch}{1.15}
\begin{tabularx}{\textwidth}{@{}>{\raggedright\arraybackslash}p{2.1cm}>{\raggedright\arraybackslash}p{4.6cm}>{\raggedright\arraybackslash}Xcc@{}}
\toprule
Tier & Simplest sufficient evidence & Example (dev question) & Dev & Test \\
\midrule
\textbf{T1}\newline local fact & one coherent passage, list or table region gives the answer and the labels that bind it to its subject and constraints & \emph{Whose consent is needed before a prosecution under the Bribery Act 2010 can be brought in England and Wales?} (LAW-en) & 18 & 30 \\
\addlinespace
\textbf{T2}\newline within-document integration & one document suffices, but no single region does; definition bridge, rule plus exception, comparison or derivation, bounded synthesis & \emph{Wise's FY2024 annual report gives its total customer balances and, in the Wise Business section, the business deposits it holds. Roughly what share of total customer balances were business deposits?} (RPT-en, derivation) & 18 & 30 \\
\addlinespace
\textbf{T3}\newline cross-document integration & every sufficient bundle spans at least two documents; cross-time or cross-version comparison, cross-subject comparison, cross-document bridge, aggregation & \emph{Turkcell and Vodafone Turkey both had their 2G licences extended to 30 April 2029. What fee did each operator agree for the extension, and how did the way they paid it differ?} (RPT-en, cross-subject) & 18 & 30 \\
\addlinespace
Unanswerable & no sufficient evidence in the pool (absence audit) & \emph{How many members did Zalando Plus have at the end of 2024?} (RPT-en) & 6 & 6 \\
\bottomrule
\end{tabularx}
\caption{The three question tiers. A tier is set by the simplest sufficient evidence in the pool; for T2 and T3, removing any required region makes part of the answer impossible. Each pool has 3 dev and 5 test questions per tier plus one unanswerable question per split. Figure~\ref{fig:example} (Appendix~\ref{app:benchmark}) shows a Persian T3 question.}
\label{tab:tiers}
\end{table*}

% src: work/multidoc/pools/<pool>.json (doc counts, summed chars); work/datasets_v2/questions/*_{dev,test}.jsonl (question counts)
% R4: sources moved here from the benchmark prose (same facts: 97 reports / 42 Iranian companies; 15 reports / 11 companies; named works)
\begin{table*}[t]
\centering
\small
\begin{tabularx}{\textwidth}{@{}lrrrr>{\raggedright\arraybackslash}X@{}}
\toprule
Pool & Docs & Chars (M) & Dev & Test & Sources \\
\midrule
RPT-fa (reports) & 97 & 4.76 & 10 & 16 & natural annual and activity reports of 42 Iranian companies \\
RPT-en (reports) & 15 & 16.64 & 10 & 16 & annual or integrated reports of 11 international companies \\
BOOK-fa (non-fiction) & 9 & 5.70 & 10 & 16 & modern Persian prose transcribed on Wikisource (e.g., Kasravi, Pirnia, Bahar) \\
BOOK-en (non-fiction) & 11 & 6.26 & 10 & 16 & Project Gutenberg non-fiction (e.g., Darwin, Mill, Tocqueville, Douglass) \\
LAW-fa (statutes) & 13 & 1.18 & 10 & 16 & Iranian statutes from Wikisource transcriptions (Civil Code, Commercial Codes, Direct Taxes Acts and others) \\
LAW-en (statutes) & 10 & 3.52 & 10 & 16 & EU regulations (GDPR, DSA, AI Act and others) and UK Acts \\
\midrule
Total & 155 & 38.06 & 60 & 96 & 58,218 chunks of at most 1,000 characters (mean about 810; 200 overlap; Appendix~\ref{app:implementation}) \\
\bottomrule
\end{tabularx}
\caption{Search pools, sources and question counts. Each pool is searched in full; dev and test share pools but no evidence document, source block, information need or source family.}
\label{tab:pools}
\end{table*}


The documents come from three genres, each in Persian and in English: company reports, historical or expository non-fiction, and statutes. Table~\ref{tab:pools} lists the six search pools and their sources. We used only born-digital text or verified transcriptions, never scans, and every question is answered against its whole pool, with no document picked in advance.

The tiers (Table~\ref{tab:tiers}) are defined by the simplest evidence that suffices anywhere in the pool. A T1 question asks for a local fact. A T2 question needs several regions of one document, so no single passage answers it. A T3 question needs at least two documents, so no single document answers it. Each genre and language holds the same number of questions per tier, and each split adds six unanswerable questions, one per pool. Our review budget set the split sizes: 60 questions for development (dev) and 96 for the test, of which 90 are answerable.

Many test questions share source material, so we do not treat them as independent. Our confidence intervals resample whole groups of questions that share a source (20 groups in the test; Table~\ref{tab:composition}, Appendix~\ref{app:benchmark}). Test questions mix ordinary and paraphrased wording.

Each question carries three kinds of annotation. These are the atomic requirements an answer must meet, the evidence passages with exact quotes and offsets into frozen text, and one or more \emph{evidence bundles}. A bundle is a set of passages that together suffice. A pipeline sees only the question, its language and its pool. The shared answer prompt also receives the expected answer format (single fact, list, comparison, series or explanation), identically for every pipeline and reader. Requirements, evidence and expected answers are visible only to the evaluator.

% src: tables/pools.tex (Persian pools 4.76+5.70+1.18 = 11.64M characters, English 16.64+6.26+3.52 = 26.42M); 45 answerable test questions per language (tab:tiers); within-language detectable difference about 8 of 45: S14_bilingual/RESULTS.md sec. 2.3 (8.3 of 45 at discordance 0.2)
\emph{What the two languages allow.} The Persian and English halves are not translations of each other. They use different documents, with about 12 million characters of Persian text and 26 million of English (Table~\ref{tab:pools}), and each half has 45 answerable test questions. A difference between the halves is therefore a difference between languages and sources together. The design supports three kinds of claim. The first concerns properties of a language or its tooling, such as the token cost of a character, which we measure directly. The second asks which findings appear in both halves, which is a form of replication. The third uses designs in which every question is its own control, such as asking every question in the other language. A paired comparison within one half detects only a difference of about 8 answers out of 45. We therefore treat any split by language as descriptive unless it comes from one of these designs.

The questions were built in six steps, four automated and two by people. Claude Opus 5.5 agents wrote a specification for each question and drafted the question from it, and a mechanical verifier and a separate Claude Opus 5.5 source-review agent checked them. A native Persian speaker then reviewed the Persian wording, and we spot-checked 32 of the 156 questions (31 accepted, 1 fixed). Dev and test share no evidence document, source block or information need, but they share the searchable pools. Appendix~\ref{app:benchmark} gives the schema and the construction steps.

\subsection{Pipelines}
\label{sec:systems}
Each pipeline turns a question into an \emph{evidence pack}, the retrieved text that the reader sees (Figure~\ref{fig:pipelines}, Table~\ref{tab:settings}). All four share the same chunking, answer prompts and readers. They differ in what they retrieve and how they pack it into the budget, so we compare whole pipelines and do not credit differences to single components. Sections~\ref{sec:components} and~\ref{sec:lost} vary some components afterwards.
% src: code-derived schematic, no estimated quantities. RAG-Paper-upgrades/ragbench: retrieval/hybrid.py,
%      approaches/hybrid.py (RRF depth 40; rerank: 50 fused candidates rescored, top 40 packed), approaches/lex5b/approach.py
%      (depth 40, top 24 shown, view 'literal', 3 s deadline, fallback = fusion order), approaches/lex5/bundles.py (exact
%      bundle packing + fill), approaches/lex5b/views.py, approaches/setr/approach.py (top 20, PASSAGE_CAP 1500, 8 s,
%      fallback = HYBRID order); test_freeze/PROTOCOL.md sec. 2 and 9 (readers, judge, human review GLM 6,000, 30 q).
\begin{figure*}[t]
\centering
\begin{tikzpicture}[
  x=1cm, y=1cm,
  box/.style={stage=#1, text width=2.8cm, minimum height=1.3cm},
  neutral/.style={draw=cAxis!70, line width=0.8pt, rounded corners=4pt, fill=black!4, align=center, font=\footnotesize,
    inner sep=3pt, text width=2.8cm, minimum height=1.3cm, text=cInk},
  ghost/.style={draw=cAxis!45, line width=0.8pt, dashed, rounded corners=4pt, align=center, font=\footnotesize,
    inner sep=3pt, text width=2.8cm, minimum height=1.3cm, text=cAxis},
  head/.style={font=\small\bfseries, text=cInk, anchor=south},
  tag/.style={draw=#1, fill=white, line width=0.8pt, rounded corners=2pt, font=\scriptsize\bfseries, text=#1,
    inner sep=1.6pt},
]
\def\cA{3.30}\def\cB{6.72}\def\cC{10.14}\def\cD{13.95}
% ---------------------------------------------------------------- column headers
\node[head] at (\cA,0.86) {Candidates};
\node[head] at (\cB,0.86) {System-specific step};
\node[head] at (\cC,0.86) {Evidence pack ($\le B$)};
\node[head] at (\cD,0.86) {Reader and scoring};
\draw[cGrid, line width=0.6pt] (0.1,0.8) -- (15.45,0.8);
% ---------------------------------------------------------------- lane labels
\foreach \y/\c/\n/\s in {0/cHyb/\hybrid/no extra call, -2.0/cHR/\hr/cross-encoder, -4.0/cLex/\lexb/1 LLM call, -6.0/cSetR/\setr/1 LLM call}{
  \node[lanelabel=\c] at (1.8,\y+0.14) {\n};
  \node[font=\footnotesize, text=cAxis, anchor=east] at (1.8,\y-0.26) {\s};
}
% ---------------------------------------------------------------- candidates
\node[neutral] (r1) at (\cA,0)     {BM25 top 40 + dense top 40, rank fusion (RRF)};
\node[neutral] (r2) at (\cA,-2.0) {same retrieval; up to 50 fused candidates};
\node[box=cLex] (r3) at (\cA,-4.0) {BM25 + dense + learned sparse; score fusion, 40 packages};
\node[neutral] (r4) at (\cA,-6.0) {\hybrid{} candidates: top 20 chunks, shown whole};
% ---------------------------------------------------------------- system-specific step
\node[ghost] (s1) at (\cB,0) {no selection step};
\node[box=cHR] (s2) at (\cB,-2.0) {cross-encoder rescores the 50; keeps the top 40};
\node[box=cLex] (s3) at (\cB,-4.0) {reads excerpts of the top 24 of the 40 packages; names bundles};
\node[box=cSetR] (s4) at (\cB,-6.0) {SetR prompt with chain of thought selects a set};
\node[tag=cHR]   at (s2.north west) [xshift=0.42cm, anchor=south, yshift=-0.5pt] {CE};
\node[tag=cLex]  at (s3.north west) [xshift=0.42cm, anchor=south, yshift=-0.5pt] {LLM};
\node[tag=cSetR] at (s4.north west) [xshift=0.42cm, anchor=south, yshift=-0.5pt] {LLM};
\node[tag=cLex]  at (s3.north east) [xshift=-0.62cm, anchor=south, yshift=-0.5pt] {3\,s limit};
\node[tag=cSetR] at (s4.north east) [xshift=-0.62cm, anchor=south, yshift=-0.5pt] {8\,s limit};
% ---------------------------------------------------------------- packing
\node[neutral] (p1) at (\cC,0)     {ranked chunks, fused order; last chunk clipped};
\node[neutral] (p2) at (\cC,-2.0) {ranked chunks, reranked order; last chunk clipped};
\node[box=cLex] (p3) at (\cC,-4.0) {exact bundle packing of whole packages, then fill; no clipping};
\node[neutral] (p4) at (\cC,-6.0) {selected set first, then \hybrid{} fill; last chunk clipped};
% ---------------------------------------------------------------- flows
\foreach \i/\c in {1/cHyb,2/cHR,3/cLex,4/cSetR}{
  \draw[flow=\c] (r\i) -- (s\i);
  \draw[flow=\c] (s\i) -- (p\i);
}
% fallbacks: a missed time limit or an unusable reply skips the selection
\draw[fallback=cLex] (r3.south) -- ++(0,-0.17) -| node[pos=0.25, fill=white, inner sep=1pt, font=\footnotesize, text=cLex] {fallback: fusion order} (p3.south);
\draw[fallback=cSetR] (r4.south) -- ++(0,-0.17) -| node[pos=0.25, fill=white, inner sep=1pt, font=\footnotesize, text=cSetR] {fallback: \hybrid{} order} (p4.south);
% ---------------------------------------------------------------- shared reader and scoring
\node[shared, text width=2.6cm, minimum height=2.6cm] (reader) at (\cD,-1.0)
  {\textbf{Reader}\\[1pt] one answer prompt; GLM-5.3-flash (primary); 5 readers replay the same packs};
\node[shared, text width=2.6cm, minimum height=2.6cm] (judge) at (\cD,-5.0)
  {\textbf{Scoring}\\[1pt] strict LLM judge on every answer; human review: GLM, 6,000 chars, 30 selected questions};
\coordinate (bus) at (11.95,0);
\foreach \i in {1,2,3,4}{\draw[line width=1.1pt, cAxis] (p\i.east) -- (p\i.east -| bus);}
\draw[line width=1.1pt, cAxis] (p1.east -| bus) -- (p4.east -| bus);
\draw[flow] (bus |- reader.west) -- (reader.west);
\draw[flow] (reader) -- (judge);
\end{tikzpicture}
\caption{The four pipelines. They share chunking, an answer prompt and a maximum evidence budget $B$; coloured boxes mark steps specific to one pipeline. The two LLM selectors fall back on a timeout or unusable reply. Five more readers replay the primary run's packs. Settings are in Table~\ref{tab:settings}.}
\label{fig:pipelines}
\end{figure*}

% src: RAG-Paper-upgrades/ragbench: retrieval/hybrid.py, config.yaml (chunking), approaches/hybrid.py (rerank),
%      approaches/lex5b/{approach,views}.py, approaches/lex8/fusion.py; work/datasets_v2/setr/SETR.md;
%      test_freeze/PROTOCOL.md sec. 2 and 9 (arms, readers, pins). Same values as the R3 systems prose.
\begin{table*}[t]
\centering
\small
\renewcommand{\arraystretch}{1.12}
\begin{tabularx}{\textwidth}{@{}>{\raggedright\arraybackslash}p{2.2cm}>{\raggedright\arraybackslash}p{2.6cm}>{\raggedright\arraybackslash}p{2.9cm}>{\raggedright\arraybackslash}X>{\raggedright\arraybackslash}p{3.3cm}@{}}
\toprule
 & \hybrid & \hr & \lexb & \setr \\
\midrule
Channels & BM25, dense & BM25, dense & BM25, dense, learned sparse & BM25, dense \\
Candidates & top 40 per channel; equal-weight RRF & top 50 fused, rescored; top 40 packed & top 40 per channel; three-signal score fusion & selection over \hybrid's top 20 full chunks \\
Extra step & -- & cross-encoder bge-reranker-v2-m3 (local GPU, no LLM tokens) & one selector call over excerpts of the top 24 units: at most 8 parts, at most 2 bundles per part & one chain-of-thought call with the SetR prompt \\
Deadline / cap & -- & -- & 3\,s / 600 completion tokens & 8\,s / 2,000 completion tokens \\
Packing & fused order; last chunk clipped & reranked order; last chunk clipped & best feasible bundle combination (exact cost, no clipping), then fused units that fit whole & selected set first, filled from \hybrid's top 40; clipped \\
On failure & -- & -- & fusion order (\emph{fallback pack}) & \hybrid{} order \\
\midrule
Shared & \multicolumn{4}{>{\raggedright\arraybackslash}p{13.0cm}}{1,000-character chunks (200 overlap); budgets $B\in\{6{,}000;\,15{,}500\}$ rendered characters; one answer prompt per answer format; selector model and primary reader GLM-5.3-flash; five replay readers; judge gpt-oss-120b (Table~\ref{tab:settings-full})} \\
\bottomrule
\end{tabularx}
\caption{The four pipelines and their settings. Further settings, model versions and provider pins are in Appendix~\ref{app:implementation}.}
\label{tab:settings}
\end{table*}


\hybrid{} retrieves chunks with BM25 and with dense bge-m3 vectors \citep{chen2024bgem3} and merges the two rankings by reciprocal rank fusion \citep{cormack2009rrf}. \hr{} rescores the merged list with a cross-encoder reranker, a trained model that scores each passage against the question and needs no LLM tokens. Both fill the budget in ranked order and clip the last chunk if enough of it fits.

\lexb{} has four stages. It first merges BM25, dense and learned-sparse scores into one ranking, which is more sensitive to a strong single signal than rank fusion is. One LLM call then reads short excerpts of the top passages, splits the question into parts and, for each part, names a bundle of passages that together support it. A packer picks the combination of named bundles that covers the most parts within the budget and fills the rest with the next passages, always whole. If the call fails or misses its deadline, the pipeline falls back to the merged ranking.

\setr{} adapts the published SetR selector \citep{lee2025setr}. One chain-of-thought call uses SetR's own prompt to pick a set from \hybrid's candidates. The selected chunks come first and \hybrid's list fills the rest of the budget. Our version uses an untuned selector model and adds the fill, so its results say nothing about the published system.

\label{sec:negative}
\lexb{} is the last of many designs tried on dev, and most variants failed there (Table~\ref{tab:negative-full} in Appendix~\ref{app:negative}). Two findings shaped the final design. First, the three candidate lists together already held a complete gold bundle for 48 of 54 answerable dev questions. Yet the packs of the non-LLM pipelines delivered far fewer, so on dev evidence was lost between candidate list and pack. On the test, more evidence is lost before the candidate list (Section~\ref{sec:lost}). Second, extra query-time lexical search added almost nothing, index-time enrichment did not transfer, and pruning the pack to the selected bundles alone lost evidence the selector had missed. A selection call over good candidates, with a fill, was what remained.

\subsection{Readers and Judge}
\label{sec:readers}
One reader, GLM-5.3-flash, answers every question and also makes both selection calls. Five readers from other model families answer the identical evidence packs, so they add breadth but not independent replications (Table~\ref{tab:settings-full}). We use open models that a deployment without access to the large commercial APIs could run. A pinned LLM judge, which we call the first judge, scores every answer against a strict rubric. An answer counts as correct only if it meets every required point and contains no contradiction, wrong value, wrong subject or unsupported added claim. The judge occasionally returned a reply we could not parse. We leave these 35 of 4,608 answers unscored and report scores as ranges that cover every possible verdict for them. Judge prompts and settings are in Appendix~\ref{app:prompts}.

\subsection{Protocol}
\label{sec:protocol}
Table~\ref{tab:protocol} lists the rules fixed before the test.
% src: work/datasets_v2/test_freeze/PROTOCOL.md (sections 1-10); lex5b/LEX5B.md (token amendment);
%      TEST_RESULTS_HUMAN_SYM.md (symmetric review: 30 questions x 4 pipelines = 120 cells). Same facts as the R3 protocol prose.
\begin{table}[b]
\centering
\small
\renewcommand{\arraystretch}{1.1}
\begin{tabularx}{\columnwidth}{@{}>{\raggedright\arraybackslash}p{1.9cm}>{\raggedright\arraybackslash}X@{}}
\toprule
Element & Rule \\
\midrule
Development & 60 dev questions only; test files never read by the implementing agents (Claude Opus 5.5) \\
Freeze & code commit in the author's repository, manifest of input hashes, pipelines, readers, decision rules; runner refuses any other commit (re-pinned once, at the pre-test re-freeze) \\
Test run & 96 questions $\times$ 8 pipeline--budget cells, each run once; an answered cell is never re-run \\
\midrule
Accuracy gate & \lexb\textsubscript{6000} $\ge$ +5 of 90 against \emph{each} of \hybrid\textsubscript{6000} and \hr\textsubscript{6000} (scaled from +3/54 on dev) \\
Resource gate & \lexb\textsubscript{6000} query-time LLM tokens (selector + reader, reasoning included) $\le 2.5\times$ \hybrid\textsubscript{6000} (relaxed from $1.5\times$ during development) \\
Superiority & two-sided exact sign tests on discordant pairs, Holm-corrected over the two primary comparisons; paired bootstrap intervals \\
Descriptive & \setr, budget effects, other readers, language / genre / tier subgroups \\
\midrule
Unresolved judge cells & reported as identification bounds $a$..$b$; sign-test $p$ as the range over every assignment; intervals clustered by the 20 components of shared source families (all added after the freeze) \\
Human review & after the machine results, blinded to the method: all four 6,000-character pipelines on the 30 machine-discordant questions (120 cells) \\
\bottomrule
\end{tabularx}
\caption{Test protocol. All gates use the primary reader; the machine analysis is primary.}
\label{tab:protocol}
\end{table}

All method development used the dev questions, and Claude Opus 5.5 agents implemented and ran our methods without ever reading the test files. Before the first test answer we froze the code, its inputs and the decision rules, and each test question was run once in every pipeline and budget. The freeze is a commit in our repository; it was not registered with a third party.

Two gates and two tests were fixed in advance, all for the primary reader. An \emph{accuracy gate} required \lexb{} at 6,000 characters, which we call the tight budget, to beat each control, \hybrid{} and \hr, by at least 5 of the 90 answerable questions. A \emph{resource gate} limited its query-time LLM tokens. Two superiority tests covered the main comparisons, and everything else was descriptive. The protocol changed in a few places during development and after the test, mostly in operations, scoring and review. The most important change relaxed the token limit of the resource gate from 1.5 to 2.5 times \hybrid's tokens, before any \lexb{} result. Appendix~\ref{app:timeline} lists every change with its date and its timing relative to the results.

After the machine results were known, the author reviewed answers blinded to the pipeline. The review covers every question on which the four 6,000-character pipelines received different judge verdicts, and all four pipelines on each such question, 30 questions in all. We call the result the \emph{symmetric merge}, because all four pipelines were reviewed on every selected question (an earlier, uneven review is described in Section~\ref{sec:stability}). On those questions the human verdicts replace the judge's, and on all others the judge's verdicts stand. The merge was built after the machine results were known, so it is post hoc.

